% authority: science_draft
\documentclass[11pt]{article}

\usepackage[margin=1in]{geometry}
\usepackage{amsmath,amssymb,amsthm}
\usepackage{array}
\usepackage{enumitem}
\usepackage{hyperref}

\newcommand{\EqSrc}{\mathrm{EqSrc}}
\newcommand{\Block}{\mathrm{Block}}
\newcommand{\NotEqSrc}{\mathrm{not\_EqSrc}}
\newcommand{\RetainH}{\mathrm{RetainH}}
\newcommand{\GenH}{\mathrm{GenH}}
\newcommand{\SourceEquivalenceTypedObject}{\mathrm{SourceEquivalenceTypedObject\_v1}}

\title{Source-Equivalence Typed Object v1}
\author{The AEther-Flow Research Project}
\date{2026-07-07}

\begin{document}
\maketitle

\section{Control Status}

This artifact implements v18 P2-T03 as a draft/control science artifact. It
defines \(\SourceEquivalenceTypedObject\) as a source-side typed object for
future family-level \(\EqSrc\) theorem-or-countermodel work.

The artifact has no proof authority. It does not discharge general \(\EqSrc\),
adopt \(\RetainH\), adopt \(\GenH\), adopt a source law, edit canonical
ontology, adopt \(\mathrm{MetricData}(E)\), expand \(g_{\mathrm{eff}}\), derive
matter coupling, derive Einstein equations, promote a benchmark, issue a Gate
Chair verdict, authorize external outreach, or claim completed derivation.

\section{Prior Record-Local Theorem}

The prior record-local theorem candidate works inside one declared finite
source-record orbit. Let
\[
  E_{\mathrm{rec}} = (\mathrm{Obj}_{\mathrm{rec}},
  \mathrm{Inv}_{\mathrm{rec}}, \mathrm{Cmp}_{\mathrm{rec}},
  \mathrm{Ann}, \mathrm{Infl})
\]
denote that record-local packet. Inside that packet, explicitly supplied
source-side relabelings witness the equivalence-relation obligations only
when the identity, inverse, and composition checks are listed and source-side.

This prior result is not a family-level theorem. It is the local datum that
P2-T03 generalizes into typed obligations.

\section{Typed Source Family}

A typed source family is a tuple
\[
  \mathcal{F}_{\mathrm{src}} =
  (\mathcal{O}, \mathcal{M}, \mathcal{I}, C, \mathcal{A},
   \mathcal{P}, \mathcal{N})
\]
where each component is source-side:
\begin{enumerate}[label=(\alph*)]
\item \(\mathcal{O}\) is a declared class or bounded family fragment of source
objects.
\item \(\mathcal{M}\) is a declared class of admissible source morphisms or
relabelings between objects.
\item \(\mathcal{I}\) is an invariant ledger whose entries are source tokens,
source certificates, source annotations, and source influence exclusions.
\item \(C\) is a source-only comparison rule with codomain
\(\{\EqSrc,\NotEqSrc,\Block\}\).
\item \(\mathcal{A}\) is an annotation ledger recording source path, object id,
hash, rule scope, and dependency.
\item \(\mathcal{P}\) is an influence ledger for forbidden-channel and proxy
edges that must be blocked rather than ignored.
\item \(\mathcal{N}\) is a negative-control set on which the same rule returns
\(\NotEqSrc\) or \(\Block\).
\end{enumerate}

\section{Source Objects}

An object \(X \in \mathcal{O}\) is source-side data. It is not a target
topological space, smooth atlas, Lorentzian metric, detector protocol,
empirical calibration, stress-energy object, matter action, benchmark branch,
validator status, role label, handoff state, or commit state.

For P2-T03 the object-set status is \emph{explicit as a type}: the artifact
requires a declared \(\mathcal{O}\) before any family theorem can be attempted.
Concrete members of \(\mathcal{O}\) remain future theorem or countermodel data.

\section{Admissible Source Morphisms or Relabelings}

A morphism \(f:X\to Y\) is admissible only if it is declared in
\(\mathcal{M}\) and preserves the source-side information named by the ledger:
\[
  f^\ast \mathcal{I}(Y) = \mathcal{I}(X), \qquad
  f^\ast \mathcal{A}(Y) = \mathcal{A}(X),
\]
with the influence incidence in \(\mathcal{P}\) mapped or blocked by an
explicit source-side rule.

If a proposed morphism needs target topology, target metric, detector protocol,
proper-time normalization, stress-energy semantics, matter action, benchmark
behavior, registry metadata authority, validator status, role status, or
handoff status, then \(f\) is not admissible and \(C\) must return \(\Block\).

\section{Invariant Ledger}

The invariant ledger \(\mathcal{I}\) is the data that a future proof or
countermodel must preserve or refute. Each entry must identify:
\[
  (\mathrm{name}, \mathrm{source\_path}, \mathrm{object\_id},
  \mathrm{source\_hash}, \mathrm{scope}, \mathrm{dependency},
  \mathrm{proxy\_exclusion}).
\]
The ledger has family validity status \emph{unknown} in this draft. It is a
typed obligation, not a proven family invariant.

\section{Source-Only Comparison Rule}

The comparison rule is a partial source-side function
\[
  C:\mathcal{O}\times\mathcal{O}\to
  \{\EqSrc,\NotEqSrc,\Block\}.
\]
Define the typed relation
\[
  X\,\EqSrc_T\,Y \quad \Longleftrightarrow \quad C(X,Y)=\EqSrc
\]
under the side condition that all required object, morphism, invariant,
annotation, influence, negative-control, and no-target guard checks pass.

The rule returns \(\Block\) for missing annotations, unmapped proxy edges,
changed influence incidence, missing closure witnesses, target-side comparison,
or authority-status premises.

\section{Identity Closure}

Identity closure is supplied as a typed requirement:
\[
  \forall X\in\mathcal{O},\quad \mathrm{id}_X\in\mathcal{M},
  \quad C(X,X)=\EqSrc
\]
provided the invariant, annotation, influence, and no-target checks pass. A
future theorem must supply the actual identity witnesses for its chosen family
or bounded fragment.

\section{Inverse Closure}

Inverse closure is candidate-derived, not proven. The required obligation is:
\[
  \forall f:X\to Y\in\mathcal{M},\quad
  \exists f^{-1}:Y\to X\in\mathcal{M}
\]
with source-only preservation of \(\mathcal{I}\), \(\mathcal{A}\), and
\(\mathcal{P}\). If the inverse is missing, the missing inverse is a
countermodel slot rather than a theorem premise.

\section{Composition Closure}

Composition closure is candidate-derived, not proven. The required obligation
is:
\[
  f:X\to Y,\ g:Y\to Z \in \mathcal{M}
  \quad\Longrightarrow\quad g\circ f:X\to Z\in\mathcal{M}
\]
with source-only preservation of every ledger entry. If composition is not
supplied or derived, the failure is a typed obstruction slot.

\section{\(\RetainH\) Boundary}

\(\RetainH\) is a boundary condition only. It is required exactly when a future
packet asks whether source-record structure, certificate status, comparison
status, negative controls, or invariant status is preserved under an
\(H\)-indexed retention operation.

P2-T03 does not adopt \(\RetainH\). The typed object records
\(\RetainH\) status as \emph{required as a possible primitive} for broader
families and as \emph{not supplied} by this artifact.

\section{\(\GenH\) Boundary}

\(\GenH\) is a boundary condition only. It is required exactly when a future
packet constructs, enumerates, or quantifies over an \(H\)-generated source
family.

P2-T03 does not adopt \(\GenH\). The typed object records \(\GenH\) status as
\emph{required as a possible primitive} for generated families and as
\emph{not supplied} by this artifact.

\section{No-Target Guard}

The no-target guard is part of the typed object:
\[
  B=(b_{\mathrm{top}},b_{\mathrm{metric}},b_{\mathrm{det}},
  b_{\mathrm{stress}},b_{\mathrm{action}},b_{\mathrm{bench}},
  b_{\mathrm{authority}}).
\]
Every component must remain false for import-as-premise:
\[
  b_{\mathrm{top}}=b_{\mathrm{metric}}=b_{\mathrm{det}}=
  b_{\mathrm{stress}}=b_{\mathrm{action}}=b_{\mathrm{bench}}=
  b_{\mathrm{authority}}=\mathrm{false}.
\]
Thus target topology, target metric, detector protocol, stress-energy
semantics, matter action, benchmark behavior, validator status, registry
status, role identity, handoff state, commit state, and generated derivatives
are not theorem premises.

\section{Family-Level Burden}

The family-level burden is:
\[
  \mathrm{Burden}_{\EqSrc_T}(\mathcal{F}_{\mathrm{src}}):
  \quad \EqSrc_T \text{ is reflexive, symmetric, and transitive on }
  \mathcal{O}
\]
using only the source-side data and guards above.

P2-T03 does not discharge this burden. It makes the burden explicit enough for
P2-T04 audit, P2-T05 stress, and later P3 theorem-or-countermodel work.

\section{Countermodel Slots}

A minimal countermodel may target any one of these slots:
\begin{center}
\begin{tabular}{|p{0.22\linewidth}|p{0.62\linewidth}|}
\hline
Slot & Countermodel condition \\
\hline
Object domain & A declared family admits an object with no source-only ledger. \\
\hline
Morphism family & A required relabeling is missing or uses a target proxy. \\
\hline
Invariant ledger & A morphism changes a required invariant. \\
\hline
Comparison rule & \(C\) returns \(\EqSrc\) from benchmark behavior or authority status. \\
\hline
Identity & Some \(X\) lacks \(\mathrm{id}_X\). \\
\hline
Inverse & Some admissible \(f\) lacks a source-side inverse. \\
\hline
Composition & Some \(g\circ f\) is not admissible or not source-side. \\
\hline
\(\RetainH\) & Preservation under \(H\)-retention is required but not supplied. \\
\hline
\(\GenH\) & Generated-family quantification is required but not supplied. \\
\hline
No-target guard & A hidden target or authority premise is used. \\
\hline
\end{tabular}
\end{center}

\section{Distance-to-GR Effect}

Distance-to-GR effect is \emph{no distance delta}. The ledger row for
source-equivalence remains draft/control with general \(\EqSrc\), \(\RetainH\),
and \(\GenH\) open. This artifact is a new mathematical payload, not a
protected ledger update.

\section{Forbidden Conclusions}

This artifact forbids the following overreads:
\begin{itemize}
\item P2-T03 as general \(\EqSrc\) discharge.
\item P2-T03 as \(\RetainH\) adoption.
\item P2-T03 as \(\GenH\) adoption.
\item P2-T03 as source-law adoption.
\item P2-T03 as canonical ontology edit.
\item P2-T03 as detector semantics, matter semantics, coupling-law, or
matter-coupling adoption.
\item P2-T03 as stress-energy tensor, matter action, Einstein-equation, or
effective-metric authority.
\item P2-T03 as benchmark promotion, Gate Chair verdict, external outreach, or
completed derivation.
\item P2-T03 registry metadata as proof authority.
\item P2-T03 validator status as theorem premise.
\end{itemize}

\section{Source Materials}

The AEther-Flow Research Project. (2026). \emph{Recommendations implementation
plan continue task v18} [Internal implementation plan].

The AEther-Flow Research Project. (2026). \emph{Source-equivalence typed-object
problem statement v1} [Internal science draft].

The AEther-Flow Research Project. (2026). \emph{Source-equivalence typed-object
schema v1} [Internal control schema].

The AEther-Flow Research Project. (2026). \emph{Selected upstream equivalence
attempt v1} [Internal science draft].

\end{document}
